\Subhead{4.8.2}{Отрицательное дифференциальное сопротивление. Генераторы на полевых транзисторах, туннельных диодах, диодах Ганна и лавинно-пролётных диодах}

Малая переменная составляющая тока активного элемента может отдавать энергию, если реальная часть его дифференциальной проводимости отрицательна. Для параллельного контура условие старта $G_{\rm active}+G_{\rm loss}<0$, а реактивные части должны компенсироваться. В установившемся режиме нелинейность восстанавливает нулевой средний баланс; устойчивость требует уменьшения отдачи энергии при росте амплитуды.

\begin{itemize}
\item \textbf{Полевой транзистор:} источник питания отдаёт энергию через управляемый канал. Обратная связь преобразует усилительный элемент в генератор и может создать отрицательное входное сопротивление на выбранной частоте; статическая отрицательная ветвь ВАХ не обязательна.
\item \textbf{Туннельный диод:} при увеличении смещения перекрытие занятых и свободных состояний по обе стороны сильно легированного перехода может уменьшаться; ток туннелирования падает и $\dd I/\dd V<0$.
\item \textbf{Диод Ганна:} сильное поле переносит электроны в долину с меньшей подвижностью, создавая отрицательную дифференциальную подвижность. Домены сильного поля могут проходить образец; частота связана с временем пролёта и резонатором. Это объёмный эффект, а не свойство p--n-перехода.
\item \textbf{Лавинно-пролётный диод:} лавинное размножение запаздывает относительно напряжения, а пролёт носителей добавляет фазовый сдвиг. Совместно они позволяют переменному току отдавать энергию полю. Сильные лавинные флуктуации повышают шум; частота определяется динамикой и геометрией, а не одной статической ВАХ.
\end{itemize}
Отрицательное дифференциальное сопротивление — локальное свойство рабочей точки или динамического отклика. Оно не означает отрицательной полной потребляемой мощности от постоянного источника.
